\documentclass[11pt,a4paper]{article}

\def\courseTitle{Industrial Security (Bachelor)}
\def\sectionNumber{05}
\def\sectionTitle{Threat Landscape}
\def\sheetKind{Solution Sheet}

% =====================================================================
% Shared preamble for exercise sheets and solution sheets.
%
% Define BEFORE \input{}-ing this file:
%   \def\courseTitle{...}
%   \def\sectionNumber{...}
%   \def\sectionTitle{...}
%   \def\sheetKind{Exercise Sheet}   or  Solution Sheet
%
% Optional:
%   \def\courseAuthor{...}
%
% The caller must issue \documentclass{...} (typically article) BEFORE
% loading this preamble.
% =====================================================================

\usepackage[utf8]{inputenc}
\usepackage[T1]{fontenc}
\usepackage[english]{babel}
\usepackage[a4paper, margin=2.2cm]{geometry}
\usepackage{graphicx}
\usepackage{listings}
\usepackage{xcolor}
\usepackage{tikz}
\usepackage{booktabs}
\usepackage{array}
\usepackage{enumitem}
\usepackage{titlesec}
\usepackage{fancyhdr}
\usepackage{hyperref}
\usepackage{amsmath}
\usepackage{amssymb}
\usepackage{tcolorbox}
\tcbuselibrary{skins, breakable}
\usepackage{needspace}

% Discourage solo lines split across pages.
\widowpenalty=10000
\clubpenalty=10000
\raggedbottom

% =====================================================================
% Shared TikZ styles for the Industrial Security lecture series.
% Loaded by every slide deck and exercise sheet that uses TikZ.
% =====================================================================

\usetikzlibrary{
  arrows.meta,
  shapes,
  shapes.geometric,
  shapes.symbols,
  shapes.multipart,
  positioning,
  fit,
  calc,
  backgrounds,
  decorations.pathmorphing,
  decorations.pathreplacing,
  decorations.text,
  patterns,
  matrix,
  shadows
}

% --- colour palette --------------------------------------------------
\definecolor{isOT}{HTML}{1F4E79}     % OT (deep blue)
\definecolor{isIT}{HTML}{6F2DA8}     % IT (purple)
\definecolor{isSafety}{HTML}{C00000} % safety red
\definecolor{isNet}{HTML}{2E7D32}    % network green
\definecolor{isAttack}{HTML}{B23A48} % attack red
\definecolor{isAccent}{HTML}{E08E0B} % amber
\definecolor{isMute}{HTML}{6B6B6B}   % grey
\definecolor{isLight}{HTML}{EDF2F8}  % light background
\definecolor{isPLC}{HTML}{2C5282}
\definecolor{isHMI}{HTML}{2F855A}
\definecolor{isSCADA}{HTML}{6B46C1}

% --- generic node styles --------------------------------------------
\tikzset{
  isnode/.style={
    draw, rounded corners=2pt, minimum height=8mm, minimum width=18mm,
    font=\small, align=center, thick
  },
  isbox/.style={isnode, fill=isLight},
  plc/.style={isnode, fill=isPLC!15, draw=isPLC, thick},
  hmi/.style={isnode, fill=isHMI!15, draw=isHMI, thick},
  scada/.style={isnode, fill=isSCADA!15, draw=isSCADA, thick},
  sensor/.style={isnode, fill=isNet!10, draw=isNet, minimum height=6mm, minimum width=14mm, font=\scriptsize},
  actuator/.style={isnode, fill=isAccent!15, draw=isAccent, minimum height=6mm, minimum width=14mm, font=\scriptsize},
  firewall/.style={isnode, fill=isAttack!15, draw=isAttack, thick},
  server/.style={isnode, fill=isMute!15, draw=isMute!80, thick},
  switch/.style={isnode, fill=isLight, draw=black!60, minimum height=6mm, minimum width=14mm, font=\scriptsize},
  workstation/.style={isnode, fill=white, draw=isOT, thick},
  attacker/.style={isnode, fill=isAttack!20, draw=isAttack, thick, font=\small\bfseries},
  inetnode/.style={draw, ellipse, minimum width=24mm, minimum height=10mm, font=\scriptsize, fill=isLight, thick},
  process/.style={isnode, fill=isOT!15, draw=isOT, thick, minimum width=24mm},
  zone/.style={
    draw, rounded corners=4pt, thick, dashed, inner sep=4mm
  },
  conduit/.style={-{Stealth[length=2.5mm]}, thick, isNet!80!black},
  attack/.style={-{Stealth[length=2.5mm]}, thick, isAttack, dashed},
  flow/.style={-{Stealth[length=2.5mm]}, thick, isOT!70!black},
  netline/.style={thick, black!60},
  axisline/.style={-{Stealth[length=2mm]}, thick, black!70},
  tag/.style={font=\scriptsize\itshape, fill=white, inner sep=1pt},
  level/.style={
    draw, fill=isLight, minimum width=0.85\linewidth, minimum height=8mm,
    font=\small, anchor=west
  },
  brace/.style={decorate, decoration={brace, amplitude=4pt, raise=2pt}},
  legendbox/.style={draw, rounded corners=2pt, fill=isLight, inner sep=2mm, font=\scriptsize, align=left}
}

% --- handy macros ----------------------------------------------------
% \plcicon, \hmiicon etc as simple labelled boxes; useful inline.
\newcommand{\plcicon}[1][PLC]{\tikz[baseline=-0.6ex]{\node[plc, minimum width=10mm, minimum height=5mm, font=\scriptsize, inner sep=1pt]{#1};}}
\newcommand{\hmiicon}[1][HMI]{\tikz[baseline=-0.6ex]{\node[hmi, minimum width=10mm, minimum height=5mm, font=\scriptsize, inner sep=1pt]{#1};}}
\newcommand{\fwicon}[1][FW]{\tikz[baseline=-0.6ex]{\node[firewall, minimum width=10mm, minimum height=5mm, font=\scriptsize, inner sep=1pt]{#1};}}


\providecommand{\courseAuthor}{Industrial Security Lecture Series}
\providecommand{\courseInstitute}{Department of Computer Science}
\providecommand{\companyLogo}{logo}

% --- Corporate branding overrides ------------------------------------
% Picks up logo / author / brand colours from the repo-root branding.tex
% when present; falls back to the defaults above otherwise.
\IfFileExists{../../../branding.tex}{% =====================================================================
% Corporate / institutional branding -- single file to customise the
% look of the entire course (slides, notes, exercises, labs).
%
% HOW TO REBRAND IN UNDER A MINUTE:
%   1. Replace `branding/logo.pdf` with your own logo
%      (PDF preferred; PNG and JPG also work -- name it `logo.<ext>`).
%   2. (Optional) change the two colours below to your brand palette.
%   3. (Optional) override the author/institute lines below.
%   4. Re-run `make` at the repo root. That's it.
%
% Everything in this file has sensible defaults; you can delete a line
% to fall back to the built-in defaults, or comment it out.
%
% This file is loaded automatically by the slides, exercise, and lab
% preambles -- you never need to \input it manually.
% =====================================================================

% --- Logo --------------------------------------------------------------
% Path is resolved from each leaf-build directory; the default points at
% the shared `branding/` folder at the repo root. If you put your logo
% somewhere else, set the full or relative path here (without extension):
\renewcommand{\companyLogo}{../../../branding/logo}

% --- Author / institute (appears on the title page footer) ------------
\renewcommand{\courseAuthor}{}
\renewcommand{\courseInstitute}{}

% --- Brand colours ----------------------------------------------------
% slideAccent      = primary brand colour (title bands, accent rule)
% slideSecondary   = secondary brand colour (section badge, highlight)
% Both use HTML hex codes. Keep enough contrast against white text.
\definecolor{slideAccent}{HTML}{00205B}      % deep navy
\definecolor{slideSecondary}{HTML}{00A0DC}   % light blue accent
}{}

% --- listings -------------------------------------------------------
\definecolor{codebg}{rgb}{0.96,0.96,0.96}
\lstset{
  backgroundcolor=\color{codebg},
  basicstyle=\ttfamily\footnotesize,
  breaklines=true,
  frame=single,
  numbers=left,
  numberstyle=\tiny\color{black!50},
  showstringspaces=false,
  tabsize=2
}

% --- header / footer ------------------------------------------------
\pagestyle{fancy}
\fancyhf{}
\renewcommand{\headrulewidth}{0.4pt}
\renewcommand{\footrulewidth}{0pt}
\lhead{\small \courseTitle}
\rhead{\small Section \sectionNumber{} -- \sheetKind}
\cfoot{\thepage}

% --- task / solution boxes -----------------------------------------
% `before={...}` guards every task/solution box with \needspace so the
% box doesn't start with only one or two lines of breathing room at the
% bottom of a page. If less than ~9 lines remain, LaTeX bumps the box
% to the next page; otherwise it starts here and breaks normally if it
% runs long.
% Boxes are `breakable` so very long solutions don't blow the page,
% but `before={\needspace{14\baselineskip}}` pushes them to a fresh
% page whenever fewer than ~14 lines remain. That covers the vast
% majority of single-question splits in practice.
\newtcolorbox{taskbox}[1]{
  enhanced, breakable,
  colback=isLight, colframe=isOT,
  fonttitle=\bfseries\color{white},
  coltitle=white,
  colbacktitle=isOT,
  title={#1},
  arc=2pt, boxrule=0.6pt,
  left=2mm, right=2mm, top=2mm, bottom=2mm,
  before={\par\vspace{2mm}\needspace{14\baselineskip}\par},
  after={\par\vspace{2mm}}
}

\newtcolorbox{solutionbox}{
  enhanced, breakable,
  colback=isNet!5, colframe=isNet,
  fonttitle=\bfseries\color{white},
  coltitle=white,
  colbacktitle=isNet,
  title={Solution},
  arc=2pt, boxrule=0.6pt,
  left=2mm, right=2mm, top=2mm, bottom=2mm,
  before={\par\vspace{2mm}\needspace{14\baselineskip}\par},
  after={\par\vspace{2mm}}
}

\newtcolorbox{hintbox}{
  enhanced, breakable,
  colback=isAccent!10, colframe=isAccent,
  fonttitle=\bfseries\color{white}, coltitle=white, colbacktitle=isAccent,
  title={Hint},
  arc=2pt, boxrule=0.6pt
}

\newtcolorbox{warnbox}{
  enhanced, breakable,
  colback=isAttack!8, colframe=isAttack,
  fonttitle=\bfseries\color{white}, coltitle=white, colbacktitle=isAttack,
  title={Caution},
  arc=2pt, boxrule=0.6pt
}

% --- section formatting --------------------------------------------
\titleformat{\section}{\Large\bfseries\color{isOT}}{\thesection}{0.5em}{}
\titleformat{\subsection}{\large\bfseries\color{isOT!80!black}}{\thesubsection}{0.5em}{}

% --- title block macro ----------------------------------------------
\newcommand{\sheetheader}{%
  \begin{center}
    {\Large\bfseries \courseTitle}\\[0.3em]
    {\large \sheetKind{} -- Section \sectionNumber}\\[0.2em]
    {\large \sectionTitle}\\[0.2em]
    \rule{\linewidth}{0.4pt}
  \end{center}
  \vspace{0.5em}
}


\begin{document}
\sheetheader

\noindent
\textbf{Note.} Model answers are deliberately concise (3--8 sentences). Where the question is
open-ended, the rubric hint indicates what must be present for full marks. Students may give
defensible alternatives to the technique IDs and architectural proposals below; mark on
reasoning and source-citation, not on exact ID matches.

\begin{solutionbox}
\textbf{Exercise 5.1 -- Stuxnet on ATT\&CK for ICS.}\\[2pt]
\begin{tabular}{p{4.4cm} p{4cm} p{6.0cm}}
\toprule
\textbf{Phase} & \textbf{Tactic} & \textbf{Technique~ID} \\
\midrule
USB across air gap & Initial Access      & T0847 Replication Through Removable Media \\
Print-spooler 0-day spread & Lateral Movement & T0866 Exploitation of Remote Services \\
Signed-driver load & Evasion        & T0859 Valid Accounts \emph{(or} T0857 System Firmware \emph{equivalent on Windows kernel)} \\
Step7 ladder injection & Execution      & T0843 Program Download \\
VFD over-/under-speed  & Impair Process Control & T0836 Modify Parameter \\
21-second HMI replay   & Inhibit Response Function & T0856 Spoof Reporting Message \\
\bottomrule
\end{tabular}\\[3pt]
\emph{Rubric hint:} accept any technique ID that is in the right tactic column and supported
by Symantec's \emph{W32.Stuxnet Dossier} or Langner's \emph{To Kill a Centrifuge}. Penalise
answers that confuse tactics (e.g.\ calling VFD sabotage ``Impact'' rather than ``Impair
Process Control'').
\end{solutionbox}

\begin{solutionbox}
\textbf{Exercise 5.2 -- Ukraine 2015 vs.\ Ukraine 2016.}
\begin{tabular}{p{2.6cm} p{5.7cm} p{5.7cm}}
\toprule
& \textbf{2015 BlackEnergy} & \textbf{2016 Industroyer} \\
\midrule
Target tier & Three distribution oblenergos (sub-transmission / distribution). &
              Pivnichna 330\,kV transmission substation. \\
Human-in-loop & Yes -- operators saw remote attackers move the mouse and click breakers
              open via legitimate HMI sessions. &
              No -- malware framework spoke IEC 60870-5-101/104, IEC 61850 MMS and OPC DA
              directly. \\
Automation / protocols & Generic remote desktop / legitimate SCADA HMI access. &
              Modular, protocol-aware payloads + a SIPROTEC DoS (CVE-2015-5374). \\
Recovery & 1--6\,h outage for $\sim$225\,000 customers; weeks of degraded SCADA, manual
              switching, serial-to-Ethernet converters re-flashed or replaced; KillDisk on
              workstations; TDoS against call centres delayed restoration. &
              About 1\,h, $\sim$200\,MW; smaller in scope but caused by a re-usable
              \emph{framework}. \\
\bottomrule
\end{tabular}\\[3pt]
The 2015~$\to$~2016 transition shows the adversary moving from \emph{hands-on-keyboard}
abuse of legitimate tools to \emph{automated, protocol-native} tooling -- a step change in
capability and a template for later tools like Industroyer2 (2022).\\[2pt]
\emph{Sources:} SANS/E-ISAC, 2016; ICS-CERT IR-ALERT-H-16-056-01; ESET 2017; Dragos
\emph{CRASHOVERRIDE}, 2017.
\end{solutionbox}

\begin{solutionbox}
\textbf{Exercise 5.3 -- TRITON: a partial defender win.}\\
(1)~The attacker likely intended to \emph{silently neutralise} SIS trip set-points (or insert
logic that ignores trip conditions) so that a later, separate manipulation of the BPCS -- e.g.\
runaway temperature, over-pressure -- would not be caught by the last-resort safety layer,
producing a release, fire, or explosion.\\
(2)~The actual outcome was a partial defender win because a payload logic error caused the
Triconex to fail-safe trip, shutting the plant down and forcing investigation that
ultimately surfaced \texttt{trilog.exe} / \texttt{library.zip} and led to attribution
(Mandiant, 2018; OFAC sanctions of TsNIIKhM, 2020). It remains deeply unsettling because the
attacker reached \emph{into} the SIS at all -- a long-assumed red line -- and the framework was
later found at a second victim, so this was not a one-off.\\
(3)~The implant lived in the Triconex controller's RAM and spoke the proprietary
\textbf{TriStation} protocol over UDP/1502; few defenders log or parse TriStation, and most
EDR / IDS rules know nothing about Triconex internals.\\[2pt]
\emph{Rubric hint:} full marks require naming the SIS / BPCS layering and citing
Dragos~2017 or Mandiant~2017/2018.
\end{solutionbox}

\begin{solutionbox}
\textbf{Exercise 5.4 -- Volt Typhoon and LOLBins.}\\
Four LOLBins associated with Volt Typhoon (per AA24-038A and Microsoft, May~2023):
\begin{itemize}\setlength{\itemsep}{0pt}
  \item \texttt{wmic} -- WMI command-line, normally for inventory and remote queries.
  \item \texttt{ntdsutil} -- AD database tool, normally for DC maintenance / backups
        (abused to dump \texttt{NTDS.dit}).
  \item \texttt{PowerShell} -- general scripting / automation; abused for in-memory execution.
  \item \texttt{netsh} -- network configuration; abused for port-proxy / persistence
        (alternatives: \texttt{certutil}, \texttt{cmd.exe}, \texttt{dnscmd}).
\end{itemize}
Signature-based EDR keys on \emph{what runs} (file hash, known malware string, known IOC);
LOLBins are \emph{signed, on-disk, expected} Windows binaries used in expected combinations,
so there is nothing to match. The detection class that \emph{can} catch this is
\emph{behavioural / analytics}: parent--child process trees (\texttt{w3wp.exe}~$\to$~\texttt{cmd}~$\to$~\texttt{ntdsutil}
is anomalous), command-line argument analysis, baselining of admin-tool usage per host,
and identity-based detections (Kerberoasting, anomalous service-account logons). UEBA on
authentication and lateral-movement patterns closes the remaining gap.\\[2pt]
\emph{Rubric hint:} accept any defensible LOLBin from the advisory; the key learning is that
signatures fail by construction and behaviour wins.
\end{solutionbox}

\begin{solutionbox}
\textbf{Exercise 5.5 -- PIPEDREAM caught before impact.}\\
(1)~Pre-impact discovery matters geopolitically because there is no clean attribution moment
and no destroyed asset to point to -- both attacker and defender know the tool exists, but no
state needs to publicly absorb a hit or escalate in response. It also gives the entire ICS
ecosystem (vendors, asset owners, sectoral CERTs) time to harden \emph{before} the weapon
is used, which is unprecedented for an ICS-grade implant.\\
(2)~Pre-built modules covering Schneider M340/M580, Omron NX/NJ, OPC~UA and CODESYS, with
scanning, parameter changes, logic upload and denial of control, show that a state-sponsored
actor has invested in a \emph{cross-vendor platform} -- protocol expertise per device family,
modular C2, and an attack lifecycle (recon $\to$ logic upload $\to$ denial) that mirrors
mature offensive tradecraft. This is no longer a bespoke weapon like Stuxnet; it is a
re-usable kit.\\
(3)~Immediate defensive action: enforce the AA22-103A mitigations -- isolate Schneider /
Omron / OPC~UA endpoints from any routed network, change default and shared engineering
passwords, enable multi-factor authentication on engineering workstations, and deploy
protocol-aware monitoring on the OT segments hosting those devices.\\[2pt]
\emph{Sources:} CISA AA22-103A; Dragos PIPEDREAM/CHERNOVITE 2022; Mandiant INCONTROLLER 2022.
\end{solutionbox}

\begin{solutionbox}
\textbf{Exercise 5.6 -- Colonial Pipeline.}\\
Dependency chain: DarkSide encrypted the corporate IT environment, including the
\emph{billing and dispatch} systems. Colonial cannot move product it cannot bill or
schedule against -- volumetric metering, custody-transfer accounting, and the dispatch /
nomination function are run in IT, not OT -- so the operator halted the pipeline as a
business decision, not because the OT controls were down. The five-day outage was driven
by the time to restore those IT systems with confidence, not by any OT recovery work.\\[2pt]
Two architectural / organisational changes (each with a trade-off):
\begin{enumerate}\setlength{\itemsep}{0pt}
  \item \emph{Hot-standby billing / dispatch in a segmented enclave} with an offline ledger
        capable of running in degraded mode for $N$~days. Trade-off: duplicated licensing,
        ongoing data-sync and reconciliation cost, regulatory audit complexity.
  \item \emph{Manual / paper-fallback dispatch procedures} drilled regularly, with
        pre-authorised operating limits so the pipeline can keep moving product at a lower
        throughput while billing is reconciled retrospectively. Trade-off: regulator and
        customer acceptance of reduced metering precision; operational risk of running
        long with degraded telemetry.
\end{enumerate}
\emph{Rubric hint:} the IT$\to$OT dependency must be named concretely (billing /
dispatch / metering). Generic ``better backups'' is insufficient.\\
\emph{Sources:} FBI 10~May~2021 statement; GAO-22-104746.
\end{solutionbox}

\begin{solutionbox}
\textbf{Exercise 5.7 -- AA23-335A.}\\
(1)~\textbf{Unitronics Vision} series PLCs (and integrated HMIs), engineering / PCOM
protocol on TCP/20256.\\
(2)~Default password \texttt{1111}; the highest-profile victim cited is the \textbf{Municipal
Water Authority of Aliquippa, Pennsylvania} (booster station PLC defaced 25~Nov~2023, manual
fallback used, no contamination). The advisory makes clear that multiple US water and
wastewater utilities were affected by the same access pattern.\\
(3)~CISA's first mitigation is to \emph{disconnect the PLC from the public Internet} (the
advisory then lists changing the default password, enforcing MFA on remote access, and
isolating OT). Accept either ``disconnect / remove Internet exposure'' or ``change the
default password'' as the headline answer.\\
(4)~The decisive initial-access vector was \emph{Internet-exposed PLC} with default
credentials -- both the door and the key were trivial.\\[2pt]
\emph{Source:} CISA AA23-335A, 1~Dec~2023.
\end{solutionbox}

\begin{solutionbox}
\textbf{Exercise 5.8 -- Pyramid of Pain.}\\
From easiest to hardest for the attacker to change:
\begin{enumerate}\setlength{\itemsep}{0pt}
  \item \textbf{MD5 hash} -- a single recompile or one byte of padding produces a new
        hash; the attacker pays nothing.
  \item \textbf{C2 IP address} -- new VPS or proxy in minutes; minor cost.
  \item \textbf{C2 domain name} -- requires re-registration, possibly new DNS / TLS
        certificate; hours to days, modest cost.
  \item \textbf{Registry-key autorun} (host artefact) -- changing the persistence path
        means changing tooling and operator playbooks; days of work and risk of breaking
        existing implants.
  \item \textbf{ATT\&CK technique} (e.g.\ T0855) -- forces a different tradecraft; can
        mean rebuilding payloads, retraining operators, and giving up reliable access.
\end{enumerate}
\emph{Implication for SOCs:} a SOC that detects only on hashes is detecting only the cheapest
thing for the attacker to change; it will see noise and miss every targeted actor. Effective
detection must climb the pyramid -- behavioural, identity, and technique-level analytics.\\[2pt]
\emph{Source:} D.\ Bianco, \emph{The Pyramid of Pain}, 2013.
\end{solutionbox}

\begin{solutionbox}
\textbf{Exercise 5.9 -- Insider IR in 24~hours.}\\
\emph{Containment (hour~0--2).} Disable the AD / engineering accounts (do not delete --
preserve auditability), terminate active VPN sessions and Kerberos tickets, revoke
remote-access tokens, and pull badge access; recover any company-issued laptop and any
USB / external drives in possession; freeze e-mail and chat content via legal hold.\\
\emph{Evidence preservation (hour~2--8).} Snapshot the engineering server, VPN concentrator
logs, file-server access logs, EDR telemetry, DLP events, badge-system logs, and the user's
mailbox / chat history \emph{before} retention windows roll over (many SIEMs roll at 7--30
days; VPN logs often shorter). Capture forensic images of any device the engineer used in
the last 30~days using write-blockers; document chain of custody.\\
\emph{Notification (hour~0--24).} Plant manager and CISO immediately; HR and legal counsel
within hours (jurisdictional rules on employee monitoring and consent); process-safety
manager because exfiltrated tag/alarm data has process-safety implications; sectoral CERT
(e.g.\ national CSIRT) once exfiltration is confirmed; law enforcement only on legal
counsel's instruction, typically after preliminary triage. Vendor PSIRT may also be
relevant if proprietary controller config is involved.\\
\emph{Do not} confront the suspect, change the locks on their office, or wipe their laptop
before imaging it -- any of these can destroy evidence, breach employment law, or tip off
co-conspirators. Equally, do not announce internally; communications must be controlled.\\[2pt]
\emph{Rubric hint:} full marks require all four blocks (containment, evidence, notify, don't)
and at least one mention of chain of custody and one of process-safety review.
\end{solutionbox}

\begin{solutionbox}
\textbf{Exercise 5.10 -- Open Research.}\\
\emph{Rubric (5~points total).}
\begin{enumerate}\setlength{\itemsep}{0pt}
  \item \textbf{Decision named, actor named, primary source cited} (1~pt). Source must be one
        of: CISA / ICS-CERT advisory, vendor PSIRT, peer-reviewed paper, or named journalism /
        book (Greenberg \emph{Sandworm}, Zetter \emph{Countdown to Zero Day}, etc.).
  \item \textbf{Operational / commercial reasoning at the time} (1~pt). The student must
        steelman the original decision -- ``they were stupid'' is not an answer.
  \item \textbf{Alternative with its cost} (1~pt). Concrete: capex, downtime, retraining,
        regulatory exposure, contractual conflict with vendor / integrator.
  \item \textbf{Counterfactual analysis} (1~pt). Honest assessment of whether the alternative
        materially changes the outcome or only the margin -- e.g.\ MFA at Colonial \emph{prevents}
        the initial intrusion; better backups only shorten recovery.
  \item \textbf{Writing, structure, citation hygiene} (1~pt). No invented sources, no invented
        body counts, ransom amounts or attribution claims beyond cited material.
\end{enumerate}
Sample exemplar decisions students might pick (non-exhaustive): Colonial's lack of MFA on
the legacy VPN; Maersk's flat Windows domain pre-NotPetya; the choice not to firmware-update
SIPROTEC relays against CVE-2015-5374 before Industroyer; Schneider's TriStation protocol
having no authentication; default-password culture on Unitronics PLCs at US water utilities;
EternalBlue-vulnerable hosts remaining unpatched two months after MS17-010.
\end{solutionbox}

\end{document}
